\clearpage
\appendix
\section{程序与支撑材料}
\subsection{运行环境与复现命令}
\WritingContract
  {让评审能够识别正式程序入口和复现条件。}
  {列出实际软件与版本、依赖、统一运行命令、输入文件要求和输出位置。}
  {环境信息与代码清单、附件 README 及正式运行记录一致。}
  {不得记录个人绝对路径、缓存目录、临时命令或未使用依赖。}
\subsection{关键算法代码}
\WritingContract
  {展示连接数学模型与可执行实现的关键部分。}
  {按问题选取目标函数、核心约束、求解和审计代码，注明功能、输入输出及附件路径；完整可运行源码放白名单附件。}
  {代码摘录与附件源码逐字节或标记区间一致，方法、变量和参数与正文相符。}
  {不得用截图代替代码，不得粘贴全部工程文件，不得省略支撑主要结论的核心实现。}
\subsection{附件文件索引}
\WritingContract
  {用最短清单说明附件中每个正式文件的用途。}
  {列出正式源码、必要结果摘要、运行说明和审计入口，并标明对应问题。}
  {索引与最终 ZIP 白名单和代码清单一致。}
  {不得重复正文符号表、长结果表或图件，不得列历史实验、内部状态、缓存和未提交文件。}
